\begin{abstract}
\noindent
\textbf{Objetivo.} Cerrar el ciclo del Sistema de Soporte Técnico ISP (equipo ACC) con sus dos
interfaces obligatorias, su capa de calidad automatizada y su evaluación experimental frente a
ISO/IEC 25010:2023~\cite{iso25010}. \textbf{Método.} Se refactorizó el microservicio principal a
arquitectura hexagonal de cuatro capas con seis patrones GoF~\cite{gamma1994patterns},
documentados como ADR con el problema real que resolvían; se construyeron una aplicación web
(React) y una móvil (Kotlin nativo, justificada con dos criterios cuantitativos medidos sobre el
proyecto); se instrumentó el sistema con las tres señales de observabilidad de
OpenTelemetry~\cite{opentelemetry2024spec}; se implementó una pirámide de pruebas completa
(unitarias, integración, contrato con Pact, E2E en tres navegadores, carga con Locust) automatizada
en un pipeline de CI/CD de diez \emph{jobs}; y se midieron cinco características de calidad con
métricas objetivas e intervalos de confianza al 95\,\%. \textbf{Resultados.} El sistema cumple
Compatibilidad (9/9 pruebas E2E en tres navegadores) y Mantenibilidad (96.6\,\% de cobertura de
pruebas y complejidad ciclomática de 1.8, ambos dentro del umbral); Eficiencia y Fiabilidad no
alcanzan su umbral bajo
carga concurrente real, con causa raíz identificada en un límite de licencia del motor de base de
datos, no en la lógica de la aplicación; Seguridad resulta parcialmente mitigada, con una brecha
real documentada (ausencia de protección contra fuerza bruta en el inicio de sesión).
\textbf{Conclusión.} La arquitectura del sistema está diseñada para escalar y es verificable de
extremo a extremo; las brechas encontradas son reales, acotadas y accionables, no ocultadas para
mostrar una nota perfecta.
\end{abstract}

\begin{center}\textbf{Abstract}\end{center}
\begin{abstract}
\noindent
\textbf{Objective.} To close the development cycle of the ISP Technical Support System (team ACC)
with both mandatory client interfaces, an automated quality layer, and an experimental evaluation
against ISO/IEC 25010:2023~\cite{iso25010}. \textbf{Method.} The core microservice was refactored
into a four-layer hexagonal architecture with six GoF patterns~\cite{gamma1994patterns}, each
documented as an ADR against the real design problem it solved; a web client (React) and a native
mobile client (Kotlin, justified with two quantitative criteria measured on the actual project)
were built; the system was instrumented with OpenTelemetry's three observability
signals~\cite{opentelemetry2024spec}; a complete test pyramid (unit, integration, Pact contract,
three-browser E2E, Locust load) was automated in a ten-job CI/CD pipeline; and five quality
characteristics were measured with objective metrics and 95\,\% confidence intervals.
\textbf{Results.} The system meets Compatibility (9/9 E2E tests across three browsers) and
Maintainability (96.6\,\% test coverage and a cyclomatic complexity of 1.8, both within
threshold); Performance
Efficiency and Reliability fall short under real concurrent load, with root cause traced to a
database-engine license limit rather than application logic; Security is partially mitigated,
with one real documented gap (no brute-force protection on login). \textbf{Conclusion.} The
system's architecture is designed to scale and is verifiable end to end; the gaps found are real,
bounded, and actionable, not hidden to present a perfect score.
\end{abstract}
